We fix a small claim language and define, for any stream of observations, the set of \emph{admissible interpretations} and the beliefs they justify. Everything in this section is a declared benchmark assumption; the oracle never sees the world.

\paragraph{Claim language.} A domain declares entities $E$, attributes, and rules. Each attribute $a$ declares a value type, a cardinality (\texttt{single} or \texttt{multi}), whether it is \texttt{changeable} (volatile vs.\ stable), whether a report of it may be hypothesised erroneous (\texttt{error\_allowed}), and whether differing values on one key are evidence against each other (\texttt{competing\_values}; two residences compete, two nicknames coexist). A proposition is $(e,a,v)$ on key $(e,a)$. Rules are Horn clauses, \texttt{strict} or \texttt{defeasible} with declared exceptions; heads are derived attributes (e.g.\ $\mathit{work\_city}(e)\leftarrow \mathit{employer}(e,x)\wedge \mathit{hq\_city}(x,c)$ unless $\mathit{remote}(e)$).

\paragraph{Time and observations.} Valid time is an integer day; belief time is report time. An \emph{assertion} is $(\mathit{id}, t_{\mathrm{rep}}, \mathit{source}, e,a,v, \mathit{valid\_cue}, \mathit{op\_cue})$ where the valid cue is \texttt{none}, \texttt{since}~$t$ or \texttt{at}~$t$ and the operator cue is \texttt{none}, \texttt{change} (optionally \emph{from} $v_{\mathrm{prev}}$) or \texttt{correction} (\emph{of} a report id). A \emph{retraction} withdraws a report or every report of a source for all later belief times. Sources declare a reliability class (\texttt{trusted}, \texttt{standard}, \texttt{low}, \texttt{blocked}) and an \emph{origin}; sources sharing an origin are copies. A report with no temporal cue is taken to assert that the value held at least at reporting time (A1); \texttt{since}~$t$ anchors it at $t$ and \texttt{at}~$t$ asserts it at $t$. Absent later evidence of a different value, a value persists from its anchor (A4, inertia); the valid cue fixes when a value is known to hold, never when it stops, so inertia applies to \texttt{at} reports exactly as to the others.

\paragraph{Admission and withdrawal.} Policy P0: every report from a non-blocked source is admitted; reliability classes never affect admissibility, only the support layer. A correction whose source shares the origin of the corrected report \emph{withdraws} that report (A-SELF): an origin's later statement about its own earlier statement is a withdrawal, not a competing report. A correction from another origin is a competing report with a cue.

\paragraph{Dispute and contradiction.} Two admitted assertions on the same key are \emph{disputed} iff they have different values and the attribute declares competing values. A \emph{contradiction} between two accepted assertions arises only for single-valued keys when the attribute is stable, or when the attribute is changeable and the two are anchored at the same time. Dispute licenses the hypothesis of error; contradiction forbids joint acceptance. The distinction is what makes ``the London report may be wrong'' admissible even when London-then-Paris is consistent as a change.

\paragraph{Interpretations.} At belief time $\tau$, an interpretation labels every admitted assertion \textsc{true} or \textsc{err}, and every maximal run of accepted same-valued reports on a multi-valued changeable key that is followed by a different value \textsc{cont} or \textsc{ended}. It is admissible iff:
(A-ERR) a report may be \textsc{err} only if error is allowed and some \emph{accepted} competitor disputes it, or an accepted correction targets it---every error label is explained by an accepted competing report, and ``everything is wrong'' is never admissible;
(A-CORR) an accepted correction forces its target \textsc{err};
(A-CHG) an accepted change \emph{from} $v_{\mathrm{prev}}$ presupposes the previous value: if any report of $v_{\mathrm{prev}}$ is anchored strictly earlier, at least one must be accepted (a presupposition about the value, not about a particular report), and on a multi-valued key the run of $v_{\mathrm{prev}}$ preceding the change is labelled \textsc{ended} (the \emph{from} clause states that the previous value stopped);
(A-CUE, policy P0c) a report carrying a change cue is not disputed by competitors anchored strictly earlier than itself---the cue identifies the revision operator, so the earlier value is explained rather than contested---while same-time and later competitors still dispute it;
(A-CONS) no contradiction among accepted reports.
Minimality is deliberately \emph{not} imposed: preferring fewer error labels would make ``London$\to$Paris was a move'' established and ``the London report was wrong'' inadmissible, which is exactly the premature commitment the benchmark is designed to penalise. Preferences belong to the support layer. An interpretation is identified by its \textsc{err} set and the timelines it induces; keys are independent, so interpretation sets are computed per key, and the oracle is made total by a declared relaxation ladder that is counted whenever it fires (it never fires on generated streams under P0c).

\paragraph{Timelines and answers.} Under an interpretation, each accepted run holds from its first anchor; on single-valued keys it ends at an unknown change point in the gap before the next different accepted anchor; on multi-valued keys it persists (\textsc{cont}) or ends in that gap (\textsc{ended}). The value at $t$ inside a gap is $\{v_{\mathrm{prev}}\}$ or $\{v_{\mathrm{next}}\}$---unknown change dates are uncertainty by construction. Derived attributes are forward-chained per interpretation, branching over candidate sets; a defeasible rule whose exception holds contributes nothing.

\paragraph{Gold.} For a value slot, the candidates are the union over interpretations: one candidate $\Rightarrow$ \texttt{established}; several $\Rightarrow$ \texttt{unresolved} with exactly that set; none $\Rightarrow$ \texttt{unknown}. For a multi-valued slot a candidate is a \emph{set} of values and alternatives are sets (London-only, Paris-only and both are three alternatives, never flattened), and the empty set is a candidate like any other: when every interpretation justifies no member---nothing admitted on the key, or everything withdrawn---the slot is \texttt{established} with $\emptyset$, which reads ``no affiliation is justified'' and is distinct from \texttt{unknown} (no interpretation at all, or an undeclared key). A yes/no proposition true in every interpretation is \texttt{established} (false in every one: established with negative polarity---the semantics is closed over the admitted evidence, so a value never asserted does not hold); mixed $\Rightarrow$ \texttt{possible}. These three conventions (empty set, established-false, unknown) are part of the contract and, as \S\ref{sec:res-s3} reports, the one for empty sets was missing from the instruction the LLM adjudicators received. A downstream query asked right after a retraction is a \emph{warranted-retraction consequence} iff its gold is established and differs from its gold one belief step earlier; only such queries enter the primary endpoint.

\paragraph{Worked case.} With \emph{residence} declared multi-valued and competing, after ``Alex lives in London'' (o1) and ``Alex lives in Paris'' (o2) with no cues, exactly four interpretations are admissible: both accepted with London continuing (both residences), both accepted with London ended (a move), o1 in error, o2 in error. ``Did Alex move?'' is \texttt{possible}; ``Was the London report wrong?'' is \texttt{possible}; ``Which residences were reported?'' is \texttt{established} with provenance; ``Which city's tax applies?'' (a rule on residence) is \texttt{unresolved} over three alternatives and is not a warranted consequence. Had o2 said ``Alex moved from London to Paris'', the move would be established under A-CUE. Two independent implementations of this semantics (a constructive per-key enumeration and a brute-force global one) were compared on the three hand-written reference streams, where they agree on the interpretation sets and the gold of every query at every belief time; on the 500 generated streams only the weaker check that every key has at least one admissible interpretation under P0c was run (the global oracle does not scale to them). The incremental store and the symbolic replay, which share the per-key oracle, are compared on every query of every setting in \S\ref{sec:res-systems}.

\paragraph{Self-update (A-SU, variant P0cSU).} A later uncued report from the \emph{same origin} on a single-valued changeable key is, under P0c, a dispute with the origin's own earlier report. The regime analysis of \S\ref{sec:res-regime} and an external check showed this to be a poor default for self-report streams. The Bayesian argument depends on the observation model, and we state both that we considered. Let two reports of one source on one key differ in value; let $c$ be the probability that the world value changed between them and $e$ the probability that a report is erroneous, independently. Under a \emph{binary} value model (an erroneous report shows the only other value), $P(\text{later value is current}\mid\text{values differ})=\frac{c(1-e)^2+(1-c)e(1-e)}{c[(1-e)^2+e^2]+2(1-c)e(1-e)}$, which exceeds $\tfrac12$ by $c(1-2e)/(2\cdot\text{denominator})$, i.e.\ for every $c>0$ and $e<\tfrac12$ (an earlier draft omitted the $ce^2$ term, which inflated the value at large $c$). Under an \emph{open} value model (an erroneous report shows a value distinct from both the true ones), the same probability is $\frac{c(1-e)+(1-c)e(1-e)}{c+(1-c)(2e-e^2)}$, which exceeds $\tfrac12$ only when $c(1-2e)\ge(1-c)e(1-e)$, roughly when changes are at least as likely as errors. In either model the later value is the better \emph{guess} in the stated regime; a probability above one half supports a preference, not a justification, so A-SU is a declared policy choice and not a theorem. A-SU, pre-registered before it was implemented (Appendix~\ref{app:chronology}), forbids labelling a report \textsc{err} when every accepted competitor is a strictly-earlier report from its own origin and no correction targets it; it never defeats a later competitor, a correction or a withdrawal, and does not apply to stable attributes. It is adopted as the default for self-report streams and reported as the declared variant P0cSU elsewhere; its measured cost is in \S\ref{sec:res-regime}; the frozen results of Settings 1--3 are under P0c.

\paragraph{Variants.} Literal P0 (no cue shields any report: only withdrawal resolves) and P1 (class-prioritised dispute: a lower-class report disputes a higher-class one only when corroborated from a second origin) are implemented and reported as sensitivity analyses. P1 is not the main configuration because it builds a reliability ordering into the oracle that the systems under test are meant to estimate, and because it marks a trusted source's false statement \texttt{established} whenever only lower-ranked sources contradict it.
